\documentclass[12pt]{article}
\usepackage[margin=1in]{geometry}
\usepackage{amsmath, amssymb, amsthm}

\theoremstyle{definition}
\newtheorem{theorem}{Theorem}

\title{\textbf{Set Theory Foundations \& Elementary Proofs}}
\author{Rutuja Deolalikar}
\date{\today}

\begin{document}
\maketitle

\section*{Problem 1: De Morgan's Law Proof}
\begin{theorem}
For any sets $A$ and $B$, $(A \cup B)^c = A^c \cap B^c$.
\end{theorem}

\begin{proof}
We show mutual inclusion $(A \cup B)^c \subseteq A^c \cap B^c$ and $A^c \cap B^c \subseteq (A \cup B)^c$.

Let $x \in (A \cup B)^c$ be arbitrary.
\begin{align*}
    x \in (A \cup B)^c &\iff x \notin (A \cup B) \\
    &\iff \neg (x \in A \lor x \in B) \\
    &\iff (x \notin A) \land (x \notin B) \\
    &\iff (x \in A^c) \land (x \in B^c) \\
    &\iff x \in (A^c \cap B^c)
\end{align*}
Since every logical equivalence step holds bi-directionally, $(A \cup B)^c = A^c \cap B^c$.
\end{proof}

\section*{Problem 2: Matrix \& Distributive Property}
Let $A = \begin{pmatrix} 1 & 2 \\ 3 & 4 \end{pmatrix}$. The Cartesian product property holds as follows:
\begin{equation}
    A \times (B \cap C) = (A \times B) \cap (A \times C)
\end{equation}

\end{document}
